% =============================================================
% chapters/vstup.tex — ВСТУП
% Структура за п. 4 Положення КДПУ (редакція 2026-05-13):
%   актуальність → мета → завдання → об'єкт → предмет →
%   методи дослідження (з GenAI-декларацією) → практичне значення →
%   структура та обсяг роботи.
% Прибрано (PRIMARY 2026-05-13): «Матеріал дослідження» як окремий
% блок, «Апробація», «Декларація конфлікту інтересів», PRISMA з
% методів, англомовні позначки RQ/LLM, термін Specifikatsiya v2.0,
% форми ЗП(ПТ)О / П(ПТ)О → ЗПО.
% Розкриття використання генеративного ШІ — у складі «Методів
% дослідження» (НЕ окремою декларацією).
% =============================================================

\section*{\centering ВСТУП}\addcontentsline{toc}{section}{ВСТУП}

\textbf{Актуальність теми.} Професійна освіта України переживає етап системної трансформації, започаткованої Законом України «Про~професійну освіту» №~4574-IX від~12.06.2024~\cite{Zakon4574IX}, зумовленої трьома реформаційними чинниками: потребами воєнного часу (перепідготовка, евакуйовані заклади)~\cite{Kremen2024EduWar, Radkevych2022Threats}, повоєнного відновлення (ветерани, внутрішньо переміщені особи, критичні професії)~\cite{Kostrytsia2025VETRecovery, Cherniakova2025WomenVET} та~гармонізації з~Європейським Союзом (EU4Skills, ETF, EQAVET, Копенгагенський процес)~\cite{Kulachynskyi2024EU4Skills, Storonyanska2025EuVET, Ciantar2025VETPolicy, Rauner2024EuropeanVET}. За~таких обставин обласний моніторинг діяльності закладів професійної освіти (далі~-- ЗПО), передусім переміщених закладів Луганщини~\cite{NMTsPTO2025, Semenenko2023DisplacedHEI}, потребує цифрового інструментарію, який одночасно забезпечує відповідність 24~формам національної звітності за~сімома блоками, захист освітніх даних і~адаптивність до~територіальних та~функціональних ролей. Чинні національні системи (ІСУО, ЄДЕБО, ДІСО, УЦОЯО, АІКОМ) не~охоплюють потреб регіонального рівня для~професійної освіти у~повному обсязі~\cite{Bieliaieva2020InfoMech, Khozhylo2024PublicMgmt, Shatalovych2025NMT}. Водночас глобальна академічна література з~моніторингу професійної освіти 1997--2025~рр. формує теоретичні засади, які доцільно зіставити з~контекстом реформи в~Україні шляхом бібліометричного аналізу.

\textbf{Мета дослідження}~-- спроєктувати інформаційно-аналітичну систему моніторингу освітньої діяльності закладів професійної освіти Луганщини на~основі архітектури «ПрофОсвіта~Луганщини» та~реалізувати її~прототип.

\textbf{Завдання дослідження}:
\begin{enumerate}
    \item Обґрунтувати теоретичні засади моніторингу професійної освіти, проаналізувати реформу №~4574-IX та~її цифровий вимір, описати існуючі національні системи моніторингу освіти та~зіставити їх із~світовим досвідом цифрового моніторингу професійної освіти (розділ~\ref{sec:rozdil1}).
    \item Розробити концепцію та~архітектуру інформаційно-аналітичної системи «ПрофОсвіта~Луганщини» з~урахуванням специфіки переміщених закладів, цільових аудиторій, функціональних і~нефункціональних вимог; обґрунтувати архітектурні рішення (модель доступу на~основі атрибутів, цілісність аудит-журналу, ідемпотентний прийом даних) (розділ~\ref{sec:rozdil2}).
    \item Реалізувати прототип веб-системи на~базі обраного технологічного стеку, провести його~експертне оцінювання та~сформулювати дорожню карту впровадження~(розділ~\ref{sec:rozdil3}).
\end{enumerate}

\textbf{Об'єкт дослідження}~-- процес моніторингу освітньої діяльності закладів професійної освіти.

\textbf{Предмет дослідження}~-- інформаційно-аналітична система моніторингу освітньої діяльності закладів професійної освіти Луганщини.

\textbf{Методи дослідження}:
\begin{itemize}
    \item аналіз і~синтез нормативної бази (Закон №~4574-IX, Закон «Про~освіту» 2017~р., накази МОН) та~наукової літератури, типологізація;
    \item порівняльна характеристика існуючих національних систем моніторингу освіти (ІСУО, ЄДЕБО, ДІСО, УЦОЯО, АІКОМ) за~9~ключовими параметрами та~виявлення функціональної прогалини для~обласного рівня;
    \item бібліометричний аналіз~-- частотний аналіз ключових слів, побудова матриці спільної появи, кластеризація методом Лувена та~засобами VOSViewer, аналіз часової динаміки для~корпусу 1998~публікацій Web of~Science 1997--2025~рр., ансамблева інтерпретація виявлених тематичних кластерів кількома великими мовними моделями (детально~-- у~Додатку~К);
    \item проєктування інформаційної системи: реляційна модель, ABAC-схема доступу, аудит із~SHA-256~хеш-ланцюгом, ідемпотентний прийом даних;
    \item веб-розробка прототипу засобами HTML5, CSS3, JavaScript, Tailwind~CSS, Chart.js, Leaflet з~подальшим переходом на~Next.js~/~TypeScript~/~PostgreSQL у~виробничому стеку;
    \item експертне оцінювання прототипу ($n\geq 3$) із~застосуванням шкали System~Usability~Scale (SUS) та~домен-специфічних пунктів, аудит доступності за~стандартом WCAG~2.1~AA, синтетичні бенчмарки для~перевірки цілісності аудит-журналу та~часу відгуку.
\end{itemize}

Використання генеративного штучного інтелекту (моделі сімейства Anthropic Claude та~відкритих ваг) задеклароване відповідно до~Закону України «Про~академічну доброчесність» №~4742-IX і~детально описане в~Додатку~К: інструмент застосовано як~допоміжний при~написанні програмного коду, чорнових текстів і~редагуванні з~подальшою верифікацією та~авторським переписуванням; в~ансамблевому бібліометричному експерименті~-- як~один із~предметів дослідження.

\textbf{Практичне значення одержаних результатів}~-- прототип системи може~бути впроваджений у~Навчально-методичному центрі професійно-технічної освіти у~Луганській області, а~рекомендації~-- використані обласними департаментами освіти та~керівниками переміщених закладів професійної освіти Луганської області. Спроєктована архітектура та~отриманий прототип становлять основу для~подальшого масштабування системи на~інші області України.

\textbf{Структура та обсяг роботи.} Робота складається з~титульного аркуша, запевнення про~дотримання академічної доброчесності, змісту, переліку умовних скорочень, вступу, трьох~розділів, загальних висновків, списку використаних джерел~(\total{citenum}~найменувань) і~додатків. Робота містить~\totalfigures\ рисунків. Загальний обсяг основного тексту~-- \pageref{last_page}~сторінок.
